\section{Graph-Based Artificial Intelligence}
\label{section:20-graph-based-ai}

Graph-Based Artificial Intelligence (\gai) is the collection of intelligence algorithms and machine learning models that learn and reason over information represented in graph data structures.
Surveying graph processing, \ac{AI} and \ac{ML} literature shows that many algorithms and models share goals and characteristics. 
The \gai Taxonomy extends prior \gai taxonomies~\cite{Chen2020a, Smith2021, Zhou2020a, Xia2021} to include the \textit{Graph Problem Solving} and \textit{Graph Transparency} problem families, the oldest and newest work \gai respectively, and group `like' tasks together to refine prior, less-structured taxonomies~\cite{Blagec2022}\footnote{\scriptsize \url{https://openbiolink.github.io/ITOExplorer/}}.
Using our \gai taxonomy as a lens to contextualize and extend a survey of \ac{IandS} problems then shows how \textit{Graph Mining} is the most in-demand \gai task. \nrscomment{is this a relevant conclusion to the rest of your work, or just a fun fact?}

\subsection{A Taxonomy of \gai}
The \gai taxonomy results from a \ac{SLR} \nrscomment{do you really need to acronym and capitalize this one?} following an abridged version of Okoli's method~\cite{Okoli2015}. 
\nrscomment{"Our survey targeted specific \gai tasks by searching textbooks, survey papers, and scientific papers containing the term `graph' and at least one of the terms `artificial intelligence' or `machine learning'." and then cut to end of centering}
The survey target was the specific \gai tasks and was conducted by initially searching textbooks, then survey papers, and then scientific papers meeting the criteria of:
\begin{center}
    {\texttt{\scriptsize `Graph' AND (`Artificial Intelligence' OR `Machine Learning')}} \\
\end{center}
\noindent After collecting the tasks, \nrscomment{we grouped them} they were grouped thematically by goal into \gai problems and then by characteristics into a broader \ac{AI} category. 

The resulting taxonomy divides \gai into six overarching ``AI Problem'' categories (\textit{Graph Problem Solving; Graph Mining; Reasoning; Graph Abstraction; Graph Prediction} and \textit{Graph Transparency}), summarized in Table~\ref{Table:taxonomy}. 
\nrscomment{"From our survey we conclude a \gai system must account for a wide variety of possible computation patterns, ..."} The key deduction from the survey to generate the taxonomy is that the breadth of tasks in the taxonomy means that a single system trying to support all of them needs to account for varying computation patterns, dataset shapes, dataset sizes, static graphs, and dynamic graphs. 
Chapter~\ref{chapter:40-Graph-AI-Benchmarking} explores the challenges of designing a benchmark to measure performance across \nrscomment{the \gai tasks surveyed.} all \gai tasks.

\begin{table*}[t]
    \begin{center}
    \small
        \begin{tabular}{lll} 
            Relation to AI     &  AI Problem    &       Example Workloads      \\%&       \textbf{Citation(s)}    \\
            \hline
            & & \\
            Formulated as graph problem & Search                   & BFS, DFS, UC, A*, MinMax              \\%& \cite{Russell2021}            \\
                                        & Constraint Satisfaction  & Forward Checking, Backtracking        \\%& \cite{Russell2021}            \\
                                        & Planning                 & GraphPlan                             \\%& \cite{Russell2021}            \\
                                        & Scheduling               & CriticalPath                          \\%& \cite{Russell2021}            \\
                                        & Inference                & Variable Elimination (Bayes Net)     \\% & \cite{Russell2021}            \\
                                        & Reinforcement Learning   & Q-Learning                           \\% & \cite{Russell2021}            \\
            & & \\
            Operates on graph           & Pattern Matching         & Spatial Subgraph Matching             \\%& \cite{Osullivan2023}          \\
                                        & Graph Clustering         & Community Detection                  \\% & \cite{Girvan2002}             \\
                                        & Graph Partitioning       & Modularity Optimization               \\%& \cite{Brandes2006}            \\
                                        & Influence Mapping        & PageRank, BetweennessCentrality       \\%& \cite{Page1998, Brandes2001}  \\
                                        & Representation Learning  & Node2Vec, DeepWalk                    \\%& \cite{Grover2016, Perozzi2014} \\
            & & \\
            Operates on knowledge graph & Knowledge Acquisition & KG Completion, Ontology Learning   \\%& \cite{Ji2021, Maedche2004}          \\
                                        & Knowledge Representation Learning & Complex Vector Embeddings \\%& \cite{Ji2021} \\ 
                                        & Logic-Based Reasoning             & Logical Inference, Ontological Inference \\%& \cite{Chen2020} \\
                                        & Representation Based Reasoning    & Tensor Factorization, Semantic Matching  \\%& \cite{Chen2020} \\
                                        & Neural Reasoning                  & Relation Inference, Entity Resolution \\%& \cite{Chen2020} \\
            & & \\
            Graph Neural Networks       & Recurrent GNN             & Time Series Prediction               \\% & \cite{Wu2020} \\
                                        & Convolutional GNN         & Syntactic GCN, Scene-Graph Generation \\%& \cite{Zhang2019, Chang2021} \\
                                        & Graph AutoEncoders        & Link Prediction                       \\%& \cite{Wu2020} \\
                                        & Spatio-Temporal GNN       & Traffic Flow Prediction               \\%& \cite{Wu2020} \\
            & & \\
            Explainability \&  Interpretability              & Post-Hoc Explanations & XGNN                         \\% & \cite{Yuan2020} \\
                                        & Interpretable-By-Design  & GNN-Subset                           \\% & \cite{Pfiffer2022}\\

                                            
        \end{tabular}
    \normalsize
    \caption{A brief taxonomy of Graph-Based AI Problems} 
    \label{table:taxonomy}
    \end{center}
\end{table*}

\subsection{Real-World Uses of \gai}

\gai is used heavily in real-world problem-solving. 
To contextualize the \gai taxonomy, we extend a recent survey of the problems, datasets, and challenges faced by \ac{IandS} professionals published by the \ac{ARLIS}~\cite{Regli2022a}.
The survey aims to identify specific examples of \gai tasks being applied to problems with \ac{IandS} characteristics.
The survey also followed an abridged \ac{SLR} \nrscomment{acronym is a bit much} process~\cite{Okoli2015} \nrscomment{", using a set of search terms pertaining to the \ac{ARLIS} \ac{IandS} problems (`< \ac{IandS} Problem>' AND (`Artificial Intelligence' OR `Machine Learning')) and another set of terms pertaining to each task in our \gai taxonomy (`<GraphAI Task>' AND `Case Study' AND (`Intelligence' OR `Security' OR `Defense'))." and cut the rest up to "For each paper..."} and used two series of searches to identify relevant case studies. 
First, searches identify papers relevant to each of the \ac{IandS} problems identified in the \ac{ARLIS} survey using: 
\begin{center}
    \texttt{\scriptsize `<\ac{IandS} Problem>' AND (`Artificial Intelligence' OR `Machine Learning')}\\
\end{center}
\noindent Then, searching for each task from the \gai taxonomy using:
\begin{center}
    \texttt{\scriptsize `<GraphAI Task>' AND `Case Study' AND ('Intelligence' OR 'Security' OR 'Defense')}\\
\end{center}
\noindent For each paper identified as relevant, we extracted the \gai tasks and the typical input data and constraints present in the case study and grouped them by the specific \ac{IandS} domain they supported. 

The survey results are summarized in Table~\ref{Table:problem-survey}. 
The key takeaways from the survey are \nrscomment{i would make this a bulleted list so its easier to quickly identify, and because it reads awkwardly to have a comma separated list of 3 'thats'} that many of the real-world \ac{IandS} problems are network problems well suited for the application of \gai, that large dynamic knowledge graphs are a standard input format, and that many use cases have temporal constraints that preclude \gai approaches that cannot produce solutions in minutes to hours. 

\begin{landscape}
    
    \begin{table*}[t]
        \begin{center}
        \small
            \begin{tabular}{lllll} 
                AI Problem                                  & I\&S Problem                        & Input Data                       &  \gai Approach                 &   Constraints         \\
                \hline
                \textbf{\ac{GPS}}                                & Link Analysis~\cite{Xu2004}         & Knowledge Graph                   & Reachability             & Heterogeneity   \\
                %
                                                            & Navigation                          & Spatial graph                     & Path Planning            & SWaP; Complexity      \\
                %
                        ~                                   &                    ~                &                 ~                 &             ~                 &          ~            \\
                %
                \textbf{\ac{GM}}                                 & HVT Identification                  & Multi-Source Graphs               & Influence Mapping        & Heterogeneity         \\
                %
                                                            & Co-Traveller Detection~\cite{Sawas2019,Tritsarolis2021,Regli2022a,Tang2014,Wang2016a}              & Spatio-Temporal Data               & Clustering            & Temporal   \\
                %
                                                            & POI Identification~\cite{Osullivan2023,Osullivan2023a}                  & Location Data                     & Spatial Pattern Matching & Complexity \\
                
                                                            & Demographic Segmentation~\cite{Army2023,Zachary1977,Sobolevsky2014}            & Communication Logs, Social Media  & Community Detection      & Temporal     \\
                %
                                                            & Counter-Threat Finance~\cite{Keene2014,Balaskas2018,Goldsmith2020,Michalak2011,Soltani2016}              & Transaction Records, Blockchain   & Subgraph Matching        & Temporal     \\
                %
                                ~                           &                   ~                 &                 ~                 &                 ~             &          ~            \\
                %
                \textbf{\ac{GR}}                          & Kill-Chain Analysis~\cite{Tradoc2007,Hutchins2011,Sadlek2022,Osullivan2015}                & Event Logs                       & Logical Inference  & Completeness    \\
                %
                                                            & All-Source Fusion~\cite{Army2019,Regli2022a,Szekely2015}                   & Streaming Knowledge Graph         & Entity Resoultion  & Temporal   \\
                %
                                ~                           &           ~                         &                 ~                 &                 ~             &             ~         \\
                %
                \textbf{\ac{GP}}                         & Anomaly Detection~\cite{Wu2021,Ren2023}                   & Streaming Knowledge Graph         & Classification    & Temporal   \\
                %
                                                            & Link Prediction~\cite{Granovetter1973,Carmody2022,Bangerter2022, Ahmadi2023}                     & Streaming Knowledge Graph         & Completion        & Temporal   \\
                %
    

            \end{tabular}
        \normalsize
        \caption{A Survey of \gai Problems in the Intelligence \& Security Domain} 
        \label{table:problem-survey}
        \end{center}
    \end{table*}
\end{landscape}

% \subsection{\gai Challenges}
% Developing efficient and effective \gai is challenging because of the \textit{massive, heterogeneous} and increasingly \textit{dynamic} nature of input data. 
% From the problem survey, we identify a requirement to support graph processing at the billions-of-edges level on heterogeneous, dynamic graphs.
% To make the problems more tracTable, simplification approaches to the data like sampling and dimensionality reduction trade-off completeness of information to gain processing speed.
% The trade-offs become more pronounced in resource-constrained settings when time, size, weight and power are limited, users may accept a reduction in one dimension to gain in another. 
% The willingness to trade off effectiveness for efficiency for complex problems and when under externally imposed constraints is the precedent that our approach in Chapter~\ref{chapter:40-Graph-AI-Benchmarking} leverages for constructing the benchmark. 

% \subsection{\gai Way Forward}
% To begin to address the challenges outlined above, we need a mechanism to compare \gai implementations and the systems they run on.
% \textit{Benchamrks} and \textit{Workload Analyses} are the two most widely used comparison mechanisms. 

% \textbf{Benchmarks.}
% Benchmarks provide an objective mechanism to judge the relative performance of candidates on a representative problem.
% They are usually targeted to reduce uncertainty about how long a problem will take to solve for a general audience and thus focus on simple metrics.
% The focus on execution time is typically beneficial as it is usually the best indicator of a system's performance~\cite{Hennessy2019}.
% Current benchmarks for \gai programs are almost exclusively focused on simplified, static graphs and generally exclude Extract, Transform and Load (ETL) time
% Benchmarks rarely report the \textit{effectiveness} metrics along with the \textit{efficiency} metrics, making them unsuiTable for evaluating approximate or incremental approaches.  
% As a result, benchmarks provide some useful insight, but to robustly compare \gai requires an approach that goes beyond `racing' systems against each other.

% \textbf{Workload Analyses.}
% Workload analyses provide a more robust analysis of a given implementation in a particular environment on a particular workload. 
% They are typically geared towards identifying bottlenecks and informing future development to exploit or elevate those bottlenecks. 
% When workload analyses are executed on a real machine, the characterization typically requires code instrumentation and invasive system-level profiling to collect the metrics required to identify the bottlenecks. 
% As a result, most of the benefits from workload analyses come from varying the programming model, hardware allocation, or graph representation on a single system since it is difficult to compare metrics collected across systems.

% The challenge moving forward is developing a path that combines the usability of a general benchmark with the informativeness of a workload analysis without oversimplifying the analysis or needing invasive access to any system we wish to test. 




% %
% \nrscomment{instead of saying what a recipient should be able to do with it, just say a benchmark should provide xyz to enable thorough comparison and promote... tie back to improving efficiency and effectiveness in graphAI problems}
% %

% A recipient of a benchmark report should be able to identify which use cases a given combination of system, dataset, problem and programming paradigm they can expect to perform `well' in, and what tradeoffs they can make between system effectiveness and efficiency to work within their constraints. 
% This thesis provides the framework for designing that benchmark. 


% Many of the `classical' AI problems in \textit{searching, planning and scheduling} start by abstracting the complexity of the real world to a graph-based formulation before searching for solutions, planning actions, making inferences or training agents~\cite{Russell2021}.
% When the input is a graph, AI approaches focus on the graph topology or the information in the attributes of the graph. 
% Topological problems necessitate \textit{structural analysis} to generate some useful knowledge or transformation, like matching locations based on spatial similarity~\cite{Osullivan2023} or detecting latent structures in the graph~\cite{Girvan2002, Brandes2006}.
% Information problems involve \textit{reasoning} over the graph abstraction of some knowledge with programs like Cyc~\cite{Lenat1990}, FrameNet~\cite{Fillmore2003}, ConceptNet~\cite{Speer2017}, ATOMIC~\cite{Sap2019} and COMET~\cite{Bosselut2019} garnering significant attention, inching towards the idea of a general-purpose common-sense AI that has driven research for the last sixty-odd years~\cite{Mccarthy1959}.
% As machine learning grew in popularity, the desire to apply machine learning to graph data structures followed.
% The initial push developed methods for \textit{graph abstraction} to embed graphs in vector space for classical machine learning applications.
% Neural approaches have yielded techniques for generating \textit{predictions} on graphs directly, with Graph Neural Networks (GNN) for classification, link prediction and traffic analysis providing even more powerful uses of graphs. 
% Finally, the structure of graphs and the long history of their use in reasoning are gaining attention as a mechanism to improve \textit{AI Trustworthiness}.
% For example, Retrieval Augmented Generation (RAG) research is exploring how to use knowledge graphs to improve Large Language Model (LLM) output and the use of symbolic AI to ground the outputs of generative models in verified, provable facts is an emerging area of interest. 

%

% Most \gai approaches focus on analyzing the graph topology, which is suiTable for exploring search spaces in planning problems, and basic pattern matching. 
% For problems that require reasoning over the attributes of the graph, or the combination of its topology and attributes, the field still lacks appropriate approaches for dynamic graphs, heterogeneous graphs, and graphs at the exa-scale. 
% Most recently, the focus has shifted to using graphs to improve other AI models, such as with the Retrieval Augmented Generation (RAG) paradigm, where symbolic AI reasoning over graphs grounds generated content in curated, traceable facts.

% \nrscomment{replace this chunk with "We divide \gai into a taxonomy covering several overarching ``AI Problem'' categories (\textit{Searching, Planning \& Scheduling; Structural Analysis; Reasoning; Graph Abstraction; Graph Prediction} and \textit{Trustworthy AI}), which are outlined in Table~\ref{Table:taxonomy}. Each AI problem category is associated with a few key \gai problems that each involve low-level graph processing tasks. We include a few typical examples of such tasks for each \gai problem to convey the breadth of tasks that a \gai benchmarking framework must cover." then add something like "In Section xxx we discuss the challenges associated with designing a benchmark that can measure performance across all \gai tasks." or something like that to tie into why we went through all this}

% \gai problems can be broken down into several overarching ``AI Problem'' categories (\textit{Searching, Planning \& Scheduling; Structural Analysis; Reasoning; Graph Abstraction; Graph Prediction} and \textit{Trustworthy AI}) which are outlined in Table~\ref{Table:taxonomy}. 
% Each \gai problem \nrscomment{you list the tasks, the graph problems don't list them. reword.} lists several typical tasks within the \gai problem space. 
% The \gai taxonomy highlights the breadth of problems that a framework for benchmarking \gai must cover and emphasizes how difficult it is to achieve `good performance' across all \gai tasks.

% The complexities of streaming and heterogeneous graph data are discussed in detail in Section~\ref{section:32-related-benchmarking-gai}, but generally, they increase complexity and degrade the efficiency of \gai tasks. 
% Given the most commonly identified \nrscomment{by who? you or the authors of the implementation?} constraint is temporal \nrscomment{dont use parens for content that should be in your sentence, its lazy and disrupts your flow, and usually hides that the logic of the sentence is off when you drop the parens and read it directly. reading it that way, I think you never link why temporal constraints imply the tension you refer to.} (either because of the need to support real-time monitoring and reporting or because of limited data availability windows) There is a direct tension between the nature of the input data, the complexity of the \gai tasks and the constraints on processing those problems. 

% \nrscomment{I would introduce the Table instead by saying something like "IandS problems spanning xyz are graph native or otherwise require analysis of the relationships between data and can benefit from graphAI. We review several of the major ... in Table x." then maybe talk about whatever the takeaway is, I'm guessing its where the constraints are since you took the time to add that column in.}



% \nrscomment{"Most IandS problems operate on graph-native input data, with the remainder..."}
% Most input data are \textit{graph-native}, with the remainder easily represented as graph problems.

% \nrscomment{again make the IandS problems the subject of the sentence}
% Of note, almost all the input data are heterogeneous attributed multigraphs, and many of the problems are expected to receive continuous update streams. 
% Graph heterogeneity and streaming are two known sources of complexity in \gai.

% The problems are spread across four of the six domains identified in the \gai Taxonomy (Table~\ref{Table:taxonomy}).

% \nrscomment{why are we talking about these 2 problems specifically? And the point you are making with them isn't really clear. is there one?}
% \textit{Graph abstraction} is implicitly present in the \textit{graph prediction} (a downstream ML model can be used in place of a GNN), and \textit{Trustworthy AI} is an ongoing goal for all organizations, and will support human interfaces to these \gai tasks in the back-end. 
% All of the tasks are constrained by the size of the input data and the complexity of representation. 
% Of particular note, temporal limitations reflect either the requirement for near-real-time monitoring, or limited time to access data.

% Overall, the problem survey shows that \gai problems deeply proliferate in the real world. 
% \nrscomment{in real-world IandS problems}
% The complexities introduced by heterogeneous representations, temporal constraints, streaming data and the breadth of tasks further emphasize the importance of characterizing \gai performance and understanding how interactions between the representations, programs, and systems manifest and the tradeoffs that need to be made when selecting the optimal approach to solve the problem(s) facing an organization.
% \nrscomment{this says "the complexities emphasize the importance of understanding how interactions manifest and tradeoffs that must be made when choosing the approach". it's too much for one sentence. boil it down to the key point or break it up into 2. For the last bit instead of "...solve the problem(s) facing an organization" something like "... to solve IandS problems"}

% \nrscomment{Cut out the journey you took to get to the final numbers and get rid of all the hypothetical stuff. The example can be one sentence. "To classify real-time messaging behavior as spam or not for 11 million customers, a telecomm graph would have to dynamically add 548 million (nodes?) per hour(footnote that says 'Based on 8.8B msgs per day spread over 16 waking hours - cite)"}
% Consider a telecommunications provider who services one-third of a country's population. 
% Assuming that country has a population equivalent to the USA, that's about $111$ million customers.
% Customers can send messages to anyone else in the country with a phone that is willing to send and receive messages, approximately $80\%$ of the population, or around $266$ Million people. 
% The average messaging rate is around $33$ messages per person, per day\footnote{\scriptsize \url{https://intelligentcontacts.com/15-text-messaging-statistics-every-business-should-know/}}, or around $8.8$ billion text messages per day.
% Naively applying that average across the $16$ hours the average person is awake that comes out to about $548$ million messages per waking hour. 
% If, for example, that company wanted to build a graph of messages to identify individuals who might be spammers~\cite{Ren2023}, they need to be able to keep up with adding an average of $548$ million messages per hour to the graph, synthesize it with historical data and execute whatever \gai task helps them to classify a phone number as a scammer in that hour. 

% The introduction of information beyond the topological has not been studied as well. 
% \nrscomment{saying one big benchmark isnt reflective of the real world isnt particularly revealing- can you generalize a bit and say most major benchmarks?}
% Real-world graphs do not reflect these simplifying assumptions. 
%
% \nrscomment{"latent information like...." so we get a sense of why this information is important (esp in IandS)}
% It is the difference in the shapes of graphs that provide much of their utility, while also introducing complications in computation because the structure of each graph is so varied.
%
% \nrscomment{why the difference in shape provides the info isn't intuitive, you need to explain it. and clarify the second point about variability introducing complications- say something more precise. The variability makes it hard for algorithms to operate efficiently across different graph inputs or something.}
% The attributes of the vertices and edges hold huge amounts of explicitly encoded information.
% In particular, as \textit{Knowledge Graphs} (KG) become an increasingly common representation they add attribute information to the vertices and edges, complicating the topology of the graph by creating complex structures like multi-graphs and hyper-edges.
% These complications in content and structure can dramatically increase the complexity of graph processing. 

% \nrscomment{the previous bit wraps up nicely, it's a bit unexpected to then dive into each category and explain them. and then there's no wrap-up takeaway to this subsection. I almost think you don't need these, though it's nice to have a reference to what these tasks mean. I think either find a way to transition into this more smoothly, and then a way to transition out at the end, or remove this from the prose altogether and add it to a new Table of AI problem spaces. this is still the intro so the prose should move pretty quickly through all the key points you need to cover to get across what your thesis is contributing-- I don't think this cuts. I didn't read these explanations so if they end up somewhere I'll read them there.}

% \textbf{Searching, Planning and Scheduling (SPS).}
% These problems are the `classical' AI problems that are common in simple agent-based systems.
% They become \gai problems when the input data is transformed into a graph to make a simplified representation of the world and support autonomous navigation, decision making and discovery tasks in the real world. 

% \textbf{Structural Analysis (SA).}
% These problems extract latent information from the topology of graphs. 
% \gai problems in this category focus on pattern matching, discovering which data points are related and understanding how information flows through the graph. 
% In the real world, SA tasks try to match known behavioral or spatial patterns, discover subcommunities, partition resources and identify critical people in influence networks. 

% \textbf{Graph Reasoning (GR).}
% These problems are working towards building programs with common-sense~\cite{Mccarthy1959}.
% These \gai problems focus on information extraction, knowledge structuring, question answering and `thinking' based problems.
% They tend to generate knowledge graphs or accept a knowledge graph as input. 
% Reasoning tasks in the real world can include automated fact-checking, question-answering over knowledge graphs, and ontological inference tasks. 

% \textbf{Graph Abstraction (GA).}
% As machine learning gained popularity, the problem of how to encode graphs into formats usable by models emerged. 
% Initial efforts generated embeddings from graphs so that they were described by some vector of features that could be passed to a downstream machine learning model. 
% Real-world applications of graph abstraction emerge as part of processing pipelines and usually contribute to classification or prediction tasks. 

% \textbf{Graph Prediction (GP).}
% Graph prediction involves learning directly from graphs, which avoids the inherent information loss incurred by graph abstraction.
% Comprised primarily of Graph Neural Networks (GNN), these \gai tasks focus on classification, prediction and transformation problems.
% Real-world applications have used GNNs for anomaly detection, missing link prediction, and forecasting traffic flow.

% \textbf{Trustworthy AI (TAI).}
% The transparency, explainability, interpretability and trustworthiness of AI systems remain open problems in AI/ML safety~\cite{Hendrycks2021}.
% \gai tasks in this category aim to improve the trustworthiness of existing systems by combining sub-symbolic generative models with symbolic graph-based representations, providing structured knowledge to generation and making reasoning processes more transparent. 
% Using \gai to enhance trustworthiness is an emerging field of research. 

% \textbf{Handling large scale graphs.}
% While the exact details of the I\&S datasets are not available, a fictional example from an adjacent domain can illustrate how quickly \gai datasets grow. 
% To classify real-time messaging behavior as spam or not for a country roughly the size of the USA, a telecommunications service provider graph would have to dynamically add $\sim548$ million edges (message events) per hour\footnote{Based on $8.8$ billion msgs per day spread over $16$ waking hours \scriptsize ~\url{https://intelligentcontacts.com/15-text-messaging-statistics-every-business-should-know/}}~\cite{Ren2023}.
% The largest available `real world' static graph datasets are over $3.5$ billion vertices and $128$ billion edges~\cite{Dhulipala2021}, and the `huge' synthetic benchmark from Graph500 has over $4.3$ trillion vertices and $70$ trillion edges\footnote{\scriptsize \url{https://graph500.org/?page_id=12\#sec-3_4}}.
% To keep up with the growing volume and velocity of data, \gai implementations and the systems they run on must continue to improve to keep pace with expected demand.

% \textbf{Handling Heterogeneous Graphs.}
% Most of the large static graph benchmarks focus on providing large numbers of vertices and edges, edge weights, and directions for processing.
% For example, the `huge' Graph500 benchmark graph is a simple, unattributed, undirected, unweighted graph with a mostly regular structure. 
% Most research has focused on analyzing graph topology because the topological structure of vertices and edges hold huge amounts of latent information like which vertices are most influential, how information flows through the graph, what communities and subcommunities exist, and what patterns correlate.
% Knowledge Graphs complicate the structure by augmenting topologies with multigraphs and hyperedges, which are difficult to operate on efficiently. 
% Effective \gai solutions must support heterogeneous graph content organized in a variety of possible graph structures without degrading.

% \textbf{Processing Graphs Dynamically.}
% \nrscomment{"Data is generated dynamically in many domains, including those that are naturally suited to graph structures. Further, dynamic data often comes with a demand for live answers, which favors incremental approaches that produce ...(partial answers?). However, dynamic and incremental graph processing is under-explored in traditional graph algorithms and is almost completely untouched in \gai. 
% Dynamic graph processing presents a particular challenge because it requires interleaving graph construction and analysis, which means the cost of creating and updating the graph cannot be written off as a one-time, upfront penalty. To enable \gai in a dynamic setting graph representations must support fast insertion, lookup, and deletion operations, and graph algorithms must produce incremental answers ... ." }

% The velocity of big data is well established. 
% \gai applications require data to be in a graph format.
% Because it is high-velocity, data is constantly arriving and the \gai developer cannot assume that they will be operating on a static graph. 
% With new information constantly being added, updated and deleted, the required representation is a dynamic graph (streaming) graph architecture. 
% Dynamic graphs require interleaving graph construction and analysis, meaning that the costs of creating and updating the graphs cannot be assumed away as a one-time cost, driving the adoption of graph representations that have fast insertion, lookup and deletion operations.
% Further, live data correlates with a demand for live answers, and algorithms that cannot run in near-real time are unsuiTable for these applications. 
% Dynamic graph processing is under-explored in traditional graph processing and almost completely untouched concerning \gai.
% Effective \gai solutions must support dynamic graphs to remain relevant in the future. 

% \textbf{Balancing Algorithmic complexity and information loss.}

% Optimal solutions to \gai problems are broadly intracTable without complexity reduction, which comes at the cost of information loss. 
% Balancing optimality (effectiveness) with tractability (efficiency) \nrscomment{is} becomes a critical tradeoff when implementing a \gai system.
% While many of the \gai tasks within the \nrscomment{spell it out} SA problem category use heuristic approaches to \nrscomment{find} get approximate answers, very few algorithms from any category use approximate or incremental algorithms~\cite{Panyala2017, Tauer2019}.
% The tradeoff between effectiveness and efficiency may change \nrscomment{"by task and evolve" and cut till the comma} from day to day and task to task, so algorithms that generate approximate solutions initially, and refine \nrscomment{insert "them"; cut "to more optimal solutions"} to more optimal solutions over time \nrscomment{"provide more flexibility to manage the tradeoff between efficiency and effectiveness, as circumstances dictate." cut the rest from here} will give \gai developers, (and more importantly \gai users) the ability to balance these competing priorities for themselves depending on their circumstances. 


% \textbf{Operating on graphs in resource-constrained settings.}
% Processing massive heterogeneous graphs quickly enough to be useful is an expensive undertaking. 
% Budgets are limited, so adding infinitely more computing resources that increase capital and operating costs does not work economically, even for the few \gai tasks where it makes sense computationally.
% \nrscomment{replace the previous sentence with "Even when increasing compute resources is a viable path economically, it often does not make sense computationally for \gai tasks." can you tack on a brief reason it doesn't make sense for most tasks?}
% \nrscomment{"Likewise, \gai tasks are often time-constrained, since the value of streaming data decays over time, and slower insights are less likely to be impactful."} The value of streaming data decays over time, and the longer it takes to generate insights, the less impactful they are. 
% Further, \nrscomment{"storage limitations and purge requirements often prevent data from being perpetually available, so processing within a given time budget is a critical challenge."} Because of storage and other limitations, data is not perpetually available, so processing within a given time budget is an important challenge. 
% Distributing the computational burden across parallel processors relieves the computational bottleneck but introduces the \nrscomment{a} need for communication efficiency.
% Finally, the push to `move the processing to the data' by deploying \gai at-the-edge introduces Size, Weight, and Power (SWaP) constraints and \nrscomment{incentivizes} drives incremental, approximate, and other resource-preserving approaches. 
% The salient challenges to \gai are the increasing size and heterogeneous shape of inputs, the need to support dynamic graph processing, the requirement to balance effectiveness with efficiency and the omnipresent constraints of time and resources on the use of \gai to achieve these tasks in the real world.